最小费用流的每次增广都要在残量网络上找最短路，而残量网络里有反向边、费用是负的，Dijkstra 直接用不了。势能数组 $h$ 就是用来把边权「平移」成非负的。

做法是把每条边的费用换成约化费用
\[ c'(u,v)=c(u,v)+h(u)-h(v), \]
并且每轮求完最短路后令 $h(v)\leftarrow h(v)+\mathrm{dist}'(v)$，其中 $\mathrm{dist}'$ 是用约化费用算出的最短路。这样之后的约化费用依然非负——由三角不等式 $\mathrm{dist}'(v)\le \mathrm{dist}'(u)+c'(u,v)$ 整理一下，正好就是新的 $c'(u,v)\ge0$。

代价是路径长度被平移了常量 $h(s)-h(t)$，但它不影响谁是最短路。代码里每轮累加的是 $h(t)$：因为 $h(s)$ 恒为 $0$，更新后的 $h(t)$ 恰好等于原网络中 $s\to t$ 的真实最短费用。

初始势能默认全 $0$，这就要求\textbf{初始边权全部非负}；如果原图本来就有负费用边，得根据图的特性自行构建 $h$ 或者直接先跑一遍 SPFA 得到初始 $h$，再从参数传进来。
